%% Согласно ГОСТ Р 7.0.11-2011:
%% 5.3.3 В заключении диссертации излагают итоги выполненного исследования, рекомендации, перспективы дальнейшей разработки темы.
%% 9.2.3 В заключении автореферата диссертации излагают итоги данного исследования, рекомендации и перспективы дальнейшей разработки темы.
\begin{enumerate}
    \item Получены верхние оценки спектральной нормы гаусс"--~ньютоновской компоненты матрицы Гессе для полносвязных и сверточных нейронных сетей.
    \item Предложены критерии стабилизации ландшафта функции потерь при увеличении объема обучающей выборки.
    \item Для введенных критериев получены оценки скорости сходимости в окрестности точки минимума.
    \item Построены критерии достаточного размера выборки для параметрических моделей на основе стабилизации оптимизационной поверхности.
    \item Предложена теоретическая схема интерпретации связи между локальной геометрией функции потерь, архитектурой модели и законами масштабирования.
    \item Для вероятностных линейных моделей построены критерии достаточного размера выборки, основанные на стабилизации апостериорного распределения параметров.
    \item Вычислительные эксперименты качественно подтвердили согласованность полученных теоретических результатов с наблюдаемым поведением моделей.
\end{enumerate}
